\section{From quadratic models to transverse slices}\label{sec:slices}

The local model in Proposition~\ref{two:vertex} gives a vanishing statement
at the vertex of an affine quadric.  We first recover the corresponding
hypersurface test by a proper compactification.  The cone of the
hyperplane-class map is a shifted constant Tate line supported on $f=0$.
After tensoring this comparison with the ambient sheaf, the hyperplane
terms have zero nearby cycles by the smooth-chart vanishing.  Vanishing
on the whole proper quadric fiber then forces vanishing on the
hypersurface.  A second proper comparison, using projective incidence,
gives tests of arbitrary codimension.

\subsection{The projective quadric comparison}

\begin{lemma}[Quadric comparison]\label{slice:quadric}
Let $K_0$ be an algebraically closed field of characteristic different
from $\ell$, let $E=\overline{\mathbb Q}_{\ell}$, and let $B$ be a
finite-type $K_0$-scheme.  For $f\in\Gamma(B,\mathcal O_B)$ put
$i:D=V(f)\hookrightarrow B$.  If $m\geq1$, consider
\[
 g:Q_f=\left\{\sum_{a=1}^{m}u_av_a+fw^2=0\right\}
       \subset\mathbb P_B^{2m}\longrightarrow B.
\]
The powers of the hyperplane class define a distinguished triangle
\begin{equation}\label{eq:slice-quadric-triangle}
 \bigoplus_{j=0}^{2m-1}E_B(-j)[-2j]
 \longrightarrow Rg_*E_{Q_f}
 \longrightarrow i_*E_D(-m)[-2m]
 \longrightarrow
 \left(\bigoplus_{j=0}^{2m-1}E_B(-j)[-2j]\right)[1].
\end{equation}
The identification of its third term with the displayed Tate line
requires only a choice of generator of a fixed one-dimensional vector
space.  Neither $B$ nor $D$ is required to be smooth or reduced.
\end{lemma}

\begin{proof}
Write $h=c_1(\mathcal O_{Q_f}(1))\in H^2(Q_f,E(1))$.
Adjunction applied to each $h^j$ gives the displayed morphism from
$E_B(-j)[-2j]$ to $Rg_*E_{Q_f}$.  Denote the cone of their direct sum
by $\mathcal K$, and write $j_B:B\setminus D\hookrightarrow B$.

We recall the elementary cohomology calculation that controls this map.
We use the projective-space and projective-bundle calculations of
\cite[Example~16.3 and Theorem~23.2]{MilneEtale}, together with localization.
A split smooth projective quadric of dimension $n$ has no odd
cohomology, and has one copy of $E(-r)$ in degree $2r$ for
$0\leq r\leq n$, with an additional copy when $n$ is even and
$r=n/2$.  One obtains the calculation inductively by writing its
equation as $x_0x_1+q'(x_2,\ldots)=0$.  The locus $x_0\ne0$ is
$\mathbb A^n$; its complement is the projective cone over the split
quadric of dimension $n-2$.  Removing the vertex of that cone gives a
line bundle over the smaller quadric.  This yields a filtration by
affine cells, with one cell of every dimension $0,\ldots,n$ and a
second middle-dimensional cell when $n$ is even.  The initial cases
are a conic, isomorphic to $\mathbb P^1$, and the split
zero-dimensional quadric, consisting of two points.  Localization and
$R\Gamma_c(\mathbb A^r,E)=E(-r)[-2r]$ give the asserted groups.
On a smooth quadric of positive dimension, the top hyperplane power
has degree $2$, so is nonzero in $E$.  All lower hyperplane powers are
therefore nonzero as well.

Over $B\setminus D$, the geometric fibers of $g$ are smooth quadrics
of odd dimension $2m-1$.  The preceding calculation shows that the
hyperplane-power map is an isomorphism on their cohomology.  Proper
base change therefore gives $j_B^*\mathcal K=0$.

Over $D$ the family, together with its hyperplane line bundle, is the
product with
\[
 V_m=\left\{\sum_{a=1}^{m}u_av_a=0\right\}
         \subset\mathbb P_{K_0}^{2m}.
\]
Let $e=[0:\cdots:0:1]$ be its vertex.  Projection away from $e$
identifies $V_m\setminus\{e\}$ with the total space of a line bundle
over the split smooth quadric
\[
 C_m=\left\{\sum_{a=1}^{m}u_av_a=0\right\}
         \subset\mathbb P_{K_0}^{2m-1}.
\]
Compactly supported cohomology of this line bundle and localization at
the vertex give
\begin{equation}\label{eq:slice-cone-cohomology}
 H^0(V_m,E)=E,\qquad H^1(V_m,E)=0,\qquad
 H^r(V_m,E)=H^{r-2}(C_m,E)(-1)\quad(r\geq2).
\end{equation}
Thus $V_m$ has the cohomology of $\mathbb P^{2m-1}$ with one extra
copy of $E(-m)$ in degree $2m$.

The hyperplane powers remain nonzero on $V_m$.  To check the top one
without a smoothness assertion about $V_m$, resolve the vertex by the
projective line bundle
\[
 \pi:\widetilde V_m
   =\mathbb P_{\mathrm{lines},C_m}
        (\mathcal O_{C_m}(-1)\oplus\mathcal O_{C_m})
       \longrightarrow C_m.
\]
Here the subscript specifies that the projective bundle parameterizes
lines.  The inclusion of $\mathcal O_{C_m}(-1)\oplus\mathcal O_{C_m}$
into the trivial vector bundle with fiber $K_0^{2m}\oplus K_0$
defines $\widetilde V_m\to V_m$; its fiber over a point of $C_m$
parameterizes the projective line joining that point to $e$.
If $\zeta$ is the pullback
of the hyperplane class and $h_C=c_1(\mathcal O_{C_m}(1))$, the
projective bundle relation is $\zeta^2=\pi^*h_C\,\zeta$
\cite[\S23]{MilneEtale}. Hence
\[
 \int_{\widetilde V_m}\zeta^{2m-1}
       =\int_{C_m}h_C^{2m-2}=2.
\]
It follows that the top hyperplane power on $V_m$ is nonzero.
Multiplication gives the same conclusion for every lower power.
This includes $m=1$: $C_1$ consists of two points, $V_1$ consists of
two projective lines meeting at $e$, and the resolution is the
disjoint union of those lines.  The hyperplane class has degree $1$
on each component.

Consequently the cone of the vector-space morphism
\[
 \bigoplus_{j=0}^{2m-1}E(-j)[-2j]
       \longrightarrow R\Gamma(V_m,E)
\]
has exactly one nonzero cohomology group, a copy of $E(-m)$ in degree
$2m$.  Injectivity in every degree, established by the nonvanishing of
the hyperplane powers, rules out any additional cohomology group in
the cone.  The cone is therefore isomorphic to $E(-m)[-2m]$.

Proper base change identifies $i^*\mathcal K$ with the constant
pullback of this vector-space cone: both the family over $D$ and
its hyperplane-power morphism are pulled back from $V_m$.
Finally, open--closed localization gives
\[
 (j_B)_!j_B^*\mathcal K\longrightarrow\mathcal K
                  \longrightarrow i_*i^*\mathcal K\longrightarrow .
\]
Since its first term vanishes, this identifies
$\mathcal K\simeq i_*E_D(-m)[-2m]$, proving
\eqref{eq:slice-quadric-triangle}.  In particular, the third term is
constant along $D$; no middle-cohomology local system is left
unidentified. Localization and proper base change hold with finite
coefficients \cite[Tags 095L and 0DDE]{Stacks} and in their adic forms
\cite[Tags 09AH and 09C9]{Stacks}; we use these compatibilities after
inverting $\ell$ and extending coefficients to $E$.
\end{proof}

\begin{remark}\label{rem:slice-rational-coefficients}
Lemma~\ref{slice:quadric} uses rational coefficients, including when
$\ell=2$.  Its hyperplane-power comparison need not be an isomorphism
off $D$ with $\mathbb Z/2^a$ coefficients: on a smooth conic the
hyperplane class has degree $2$.  We use finite coefficients to
construct the usual cohomological compatibilities, then use
\eqref{eq:slice-quadric-triangle} after rationalization.  This imposes
no restriction on the residue characteristic.  The case $m=0$ will
be handled directly below; formally the projective model is then
$Q_f=D\subset\mathbb P_B^0$ and the hyperplane sum is empty.
\end{remark}

\subsection{The hypersurface test}

We return to the trait $S$, the bundle stack $\mathcal B$, and the
compact constructible object $F\in\mathcal D_C$ of
Proposition~\ref{spec:missing}.  For a smooth finite-type chart
$b:U\to\mathcal B$, put $F_U=b_K^*F$.  Denote by
$\mathcal N_U\subset T^*U_s$ the image of the global nilpotent cone
under smooth cotangent pullback.

\begin{proposition}[Hypersurface test]\label{slice:hypersurface}
Let $b:U\to\mathcal B$ be smooth, with $U$ a finite-type $S$-scheme.
Let $f$ be a regular function on $U$, let $D=V(f)$, and let
$u\in D_s(k)$.  If
\[
                 d(f_s)_u\notin(\mathcal N_U)_u,
\]
then
\[
                 \bigl(\Psi_D(F_U|_{D_K})\bigr)_u=0.
\]
\end{proposition}

\begin{proof}
If $d(f_s)_u$ is not in the image of $T^*_{b(u)}\mathcal B_s$, its
image in the relative cotangent space for $b$ is nonzero.  The smooth
hypersurface criterion shows that $D\to\mathcal B$ is smooth near
$u$.  Lemma~\ref{spec:vanishing} then gives the assertion.

Otherwise $d(f_s)_u=b^*A$ for a non-nilpotent covector
$A\in T^*_{b(u)}\mathcal B_s$.  Proposition~\ref{two:vertex}
provides a point $x\in\mathscr X_s(k)$ and an integer $m\geq0$ with
the following property.  Set
\[
 T=U\times_S\mathscr X,\qquad t_0=(u,x),\qquad
 D_T=D\times_S\mathscr X,
\]
and let $F_T$ be the pullback of $F_U$ to $T_K$.  For the affine
quadric over $T$
\[
 Q^a=\left\{\sum_{a=1}^m u_av_a+f=0\right\}
          \subset T\times_S\mathbb A_S^{2m},
\]
the nearby cycles of the pullback of $F_T$ vanish at
$(t_0,0,0)$.  The symbols $u_a$ here are affine coordinates, distinct
from the fixed point $u\in U_s$.
If $m=0$, this affine model is $D_T$ itself.  Smooth compatibility
for $D_T\to D$ immediately gives the desired vanishing.

Suppose $m\geq1$ and compactify to
\[
 g:Q_f=\left\{\sum_{a=1}^m u_av_a+fw^2=0\right\}
          \subset\mathbb P_T^{2m}\longrightarrow T.
\]
The fiber over $t_0$ is the split cone of
Lemma~\ref{slice:quadric}.  At every point of this fiber other than
its vertex, some $u_a$ or $v_a$ is nonzero.  In the chart $u_a=1$,
for example, its equation solves uniquely for $v_a$, so $g$ is
smooth there.  This argument works also in characteristic $2$.
Lemma~\ref{spec:vanishing} gives $\Psi_TF_T=0$, and smooth
compatibility therefore gives vanishing of the nearby cycles of
$g_K^*F_T$ at all these nonvertex points.  The vertex lies in the
chart $w=1$, where its vanishing was supplied by
Proposition~\ref{two:vertex}.

The nearby-cycle complex restricts to zero on the entire proper
fiber.  Proper compatibility and the projection formula give
\begin{equation}\label{eq:slice-quadric-pushforward-zero}
 \left(\Psi_T\bigl(F_T\otimes Rg_{K*}E\bigr)\right)_{t_0}=0.
\end{equation}
Apply Lemma~\ref{slice:quadric} on $T_K$ and tensor its triangle
with $F_T$.  The projection formula identifies the third term with
\[
             i_{K*}(F_T|_{(D_T)_K})(-m)[-2m],
             \qquad i:D_T\hookrightarrow T.
\]
After applying $\Psi_T$ and taking the stalk at $t_0$, the first
term vanishes because $\Psi_TF_T=0$, and the second vanishes by
\eqref{eq:slice-quadric-pushforward-zero}.  Proper compatibility for
the closed immersion $i$ identifies the third term with
\[
             \bigl(\Psi_{D_T}(F_T|_{(D_T)_K})\bigr)_{t_0}
                        (-m)[-2m].
\]
It too vanishes.  Smooth compatibility for $D_T\to D$ completes
the proof.
\end{proof}

The passage through the projective quadric uses no purity theorem for
the possibly singular hypersurface $D$.  It uses the actual
closed-support term of the comparison triangle.  We now use the same
principle to impose several equations simultaneously.

\subsection{Projective incidence and higher codimension}

\begin{lemma}[Incidence comparison]\label{slice:incidence}
Let $B$ be a finite-type scheme over an algebraically closed field of
characteristic different from $\ell$.  Let $f_1,\ldots,f_d$ be
regular functions on $B$, where $d\geq1$, and set
$i:J=V(f_1,\ldots,f_d)\hookrightarrow B$.  For
\[
 r:I=\left\{\sum_{a=1}^d c_af_a=0\right\}
           \subset B\times\mathbb P^{d-1}\longrightarrow B,
\]
the hyperplane powers give a distinguished triangle
\begin{equation}\label{eq:slice-incidence-triangle}
 \bigoplus_{j=0}^{d-2}E_B(-j)[-2j]
 \longrightarrow Rr_*E_I
 \longrightarrow i_*E_J(-(d-1))[-2(d-1)]
 \longrightarrow
 \left(\bigoplus_{j=0}^{d-2}E_B(-j)[-2j]\right)[1].
\end{equation}
For $d=1$ the first term is zero and $I=J$.
\end{lemma}

\begin{proof}
If $d=1$, the equation in $B\times\mathbb P^0$ is $f_1=0$,
so the stated triangle is immediate.  Suppose $d\geq2$.
Off $J$, the map $\mathcal O_B^d\to\mathcal O_B$ with coefficients
$(f_1,\ldots,f_d)$ is surjective.  Its kernel is a vector bundle of
rank $d-1$, and $I$ is its projective bundle of lines.  The powers of
the restricted hyperplane class identify its direct image with the
first term of \eqref{eq:slice-incidence-triangle}.  This can also be
checked on geometric fibers by proper base change and the cohomology
of $\mathbb P^{d-2}$.

On $J$, the family is exactly $J\times\mathbb P^{d-1}$ and its
hyperplane class comes from the second factor.  The restricted cone
of the displayed morphism is therefore the constant top class
$E_J(-(d-1))[-2(d-1)]$.  The global cone vanishes off $J$, so
open--closed localization identifies it with the pushforward of this
restricted cone, proving the triangle.
\end{proof}

\begin{proposition}[Higher-codimension test]\label{slice:higher}
Let $b:U\to\mathcal B$ be smooth, with $U$ a finite-type $S$-scheme,
and let $F_U=b_K^*F$ as above.  Suppose that $u\in U_s(k)$ and
$f_1,\ldots,f_d$ are regular functions vanishing at $u$.  Assume
that their special-fiber differentials at $u$ are linearly independent
and that
\[
 \operatorname{span}_k\{d(f_{1,s})_u,\ldots,d(f_{d,s})_u\}
                     \cap(\mathcal N_U)_u=\{0\}.
\]
Then, for $J=V(f_1,\ldots,f_d)$,
\[
                         (\Psi_J(F_U|_{J_K}))_u=0.
\]
The assertion includes $d=0$, with $J=U$.
\end{proposition}

\begin{proof}
For $d=0$ the assertion is Lemma~\ref{spec:vanishing}.
For $d\geq1$ form the proper incidence morphism
\[
 r:I=\left\{\sum_{a=1}^d c_af_a=0\right\}
       \subset U\times_S\mathbb P_S^{d-1}\longrightarrow U.
\]
Each affine projective-space chart in
$U\times_S\mathbb P_S^{d-1}$ is still smooth over $\mathcal B$.
On the chart $c_a=1$, the hypersurface equation is
$f_a+\sum_{b\ne a}c_bf_b=0$.  At a point $(u,[c])$ its
special-fiber differential is
\[
                          \sum_{b=1}^d c_b\,d(f_{b,s})_u.
\]
Indeed its projective-parameter component is zero since every
$f_b(u)$ is zero.  The displayed combination is nonzero by linear
independence, and lies outside $(\mathcal N_U)_u$ by the hypothesis.
The nilpotent cone on the product chart is the pullback of
$\mathcal N_U$ with zero parameter component.  Proposition~\ref{slice:hypersurface}
therefore gives vanishing of the nearby cycles of $r_K^*F_U$ at
every closed point over $u$.

The restriction of this nearby-cycle complex to the projective fiber
is constructible.  A nonzero constructible cohomology sheaf on a
finite-type scheme over $k$ has a nonzero stalk at a closed
$k$-point.  Thus the complex vanishes on the entire fiber.  Proper
compatibility gives
\[
                      (\Psi_U Rr_{K*}r_K^*F_U)_u=0.
\]
Tensor the triangle of Lemma~\ref{slice:incidence} on $U_K$ with
$F_U$ and apply the projection formula.  The first term is a finite
sum of shifts and twists of $F_U$, the second is
$Rr_{K*}r_K^*F_U$, and the third is
\[
                    i_{K*}(F_U|_{J_K})(-(d-1))[-2(d-1)],
                    \qquad i:J\hookrightarrow U.
\]
After specialization and taking the stalk at $u$, the first two
terms vanish.  Proper compatibility for $i$ shows that the third
is the corresponding shift and twist of
$(\Psi_J(F_U|_{J_K}))_u$, proving the result.  For $d=1$ the first
sum is empty, so the same argument remains valid.
\end{proof}

We have obtained all finite transverse-slice tests while applying the
hypersurface argument only on schemes smooth over $\mathcal B$.
The remaining task is to find one such slice on which the nonzero
constructible object $F$ has a nonzero nearby-cycle stalk.
